% docs/physics/latex/00_common.tex - SIMC++ physics note 00: the elements common to every reaction.
% Standalone PDF docs/physics/pdf/SIMCpp_Physics_00_common.pdf; Part I of the book SIMC++ Physics Reference.
% Rules of the notes: docs/physics/README.md and the header of simcnote.cls. Every display is feq/feqa with the C++
% function first and its Fortran origin (simc_gfortran upstream 59ccdf6) second. Labels: \lbl (prefix N00-); the
% reaction notes refer to them with \ceqr, \csecr, \ctabr.
\documentclass{simcnote}
\begin{document}
\simcnotemeta{00}{Common elements}{Common elements of the SIMC++ physics}%
  {The plug-in architecture, the event chain, notation and glossary, frames, kinematics, generation and weighting,
   target effects, the Coulomb correction, radiative corrections and the common cross-section ingredients}
\simcnotemaketitle
\simcabstract{What every reaction of SIMC++ shares. The first part describes how a reaction is built in SIMC++: the
stage interfaces and plug-ins, how the pipeline chooses them and how a JSON input overrides the choice, and the order
in which one event passes through them. The second part defines the notation and the terms used in all the notes, and
gives the physics common to all reactions: frames and kinematic variables, the generation limits and the event weight,
the effects of the target, the Coulomb correction, the radiative corrections, and the Mott cross section and proton
form factors used by several cross-section models. The reaction notes refer to this note instead of repeating it.}
\simcnotetoc

% =====================================================================================================================
\section{How to read the physics notes}\lbl{sec:reading}

\paragraph{One note per reaction.} Each reaction note describes one reaction as SIMC++ builds it: the reaction
plug-in and the physics-model plug-in, how they are selected, and then every formula they evaluate. The sections of a
reaction note are always the same: (1) the reaction in SIMC++, (2) kinematics, generated variables and generation
volume, (3) completing the event, (4) jacobians, (5) cross-section model, (6) radiation, energy loss and decays,
(7) inputs and outputs, (8) validation in SIMC++, (9) Fortran origin and Fortran findings. The notes are also bound as
the book \emph{SIMC++ Physics Reference}, and each reaction note is copied into the source folders of its plug-ins.

\paragraph{Citations.} Under each displayed equation two citations are printed: first the C++ function that computes
it (\csrc{File.cpp: Function}, path relative to \texttt{src/} unless stated), second the Fortran lines it was translated
from (\fsrc{file.f:lines} of SIMC, \texttt{simc\_gfortran} upstream commit \texttt{59ccdf6}). The Fortran is the
specification: SIMC++ evaluates the same formula, with the same constants and the same order of operations. Where the
code differs from the literature it cites, contains an approximation or a hard-coded number, the text says so.

\paragraph{Two modes.} By default SIMC++ corrects the clear bugs of the Fortran; the run option
\code{-{}-fortran-compat} restores the Fortran behaviour exactly. Every place where the two differ, or where an odd
Fortran behaviour is kept, is an entry of \code{docs/translation/markdown/FORTRAN\_DIFFERENCES.md}. The text of the
notes describes such places in words; their identifiers appear only in the findings table of the last section of each
note.

\paragraph{Units.} The internal SIMC units are used throughout (\tabr{tab:units}). The reaction notes state the units of
each cross section.

\begin{simctable}{small}{L{0.32\linewidth}L{0.2\linewidth}L{0.42\linewidth}}
\caption{Internal units \fsrc{constants.inc:3-9}.}\lbl{tab:units}\\
\toprule\rowcolor{simchead}\simcth{quantity} & \simcth{unit} & \simcth{remark}\\\midrule\endfirsthead
\toprule\rowcolor{simchead}\simcth{quantity} & \simcth{unit} & \simcth{remark}\\\midrule\endhead
\bottomrule\endlastfoot
energy, mass & MeV & \\
momentum & MeV/$c$ & \\
length & cm & \\
angle & rad & input angles in degrees are converted at start-up\\
slope $x'$, $y'$ & rad (dimensionless) & input acceptance limits in mrad\\
relative momentum $\delta$ & \% & \\
density & g/cm$^3$ & \\
thickness & g/cm$^2$ & input target thickness in mg/cm$^2$\\
beam charge & mC & input \code{EXPER\%charge}\\
luminosity & $\mu$b$^{-1}$ & \\
tree columns $|\vec q|$, $\nu$, $E_m$, $P_m$, $W$ & GeV, GeV/$c$ & converted when written\\
tree column $Q^2$ & (GeV/$c$)$^2$ & converted when written\\
\end{simctable}

% =====================================================================================================================
\section{The SIMC++ plug-in architecture}\lbl{sec:arch}

\subsection{Stages and plug-ins}\lbl{sec:stages}
SIMC++ builds a run from \emph{stages}. Each stage is a C++ interface (\texttt{include/simc/stages/}); each
implementation is a \emph{plug-in} registered under a name with the macro \code{SIMC\_REGISTER(Interface, "name", Class)}
in its \texttt{.cpp} file. A stage is created by name from the registry, then \code{configure(json, params)} reads its
JSON configuration and the input parameters, and \code{init(StageContext)} builds its tables and takes its
collaborators (the run set-up, the target, the radiator, the physics model). \tabr{tab:plugins} lists the plug-ins that
the physics notes refer to.

\begin{simctable}{footnotesize}{L{0.18\linewidth}L{0.17\linewidth}L{0.25\linewidth}L{0.4\linewidth}}
\caption{Stage interfaces and the built-in plug-ins used by the notes (one plug-in per row).}\lbl{tab:plugins}\\
\toprule\rowcolor{simchead}\simcth{interface} & \simcth{registered name} & \simcth{class} & \simcth{file (role)}\\\midrule\endfirsthead
\toprule\rowcolor{simchead}\simcth{interface} & \simcth{registered name} & \simcth{class} & \simcth{file (role)}\\\midrule\endhead
\bottomrule\endlastfoot
\code{IRandom} & \code{mt64} & \code{Mt64Random} & core/random/Mt64Random.cpp (default random source)\\
\code{IRandom} & \code{ranlux} & \code{RanluxRandom} & \file{core/random/RanluxRandom.cpp} (the Fortran RANLUX stream)\\
\code{ITarget} & \code{hallc\_cryo} & \code{material::HallCCryoTarget} & \file{material/HallCCryoTarget.cpp} (vertex $z$, energy loss, multiple scattering)\\
\code{IMaterialModel} & \code{simc\_material} & \code{material::SimcMaterial} & \file{material/SimcMaterial.cpp} (ionisation loss, used by the target)\\
\code{IBeam} & \code{standard} & \code{StandardBeam} & \file{reactions/StandardBeam.cpp} (beam spot, raster, beam energy)\\
\code{IReaction} & \code{hydrogen\_elastic} & \code{HydrogenElastic} & \file{reactions/HydrogenElastic.cpp} (note 01)\\
\code{IReaction} & \code{deuterium\_eep} & \code{DeuteriumEep} & \file{reactions/NuclearEep.cpp} (note 02)\\
\code{IReaction} & \code{nucleus\_eep} & \code{NuclearEep} & \file{reactions/NuclearEep.cpp} (note 02)\\
\code{IReaction} & \code{nuclear\_elastic} & \code{NuclearElastic} & \file{reactions/HydrogenElastic.cpp} (note 03)\\
\code{IReaction} & \code{phase\_space} & \code{PhaseSpace} & \file{reactions/HydrogenElastic.cpp} (acceptance runs)\\
\code{ISpectralFunction} & \code{theory} & \code{nuclear::TheorySpectralFunction} & \file{physics/nuclear/SpectralFunctions.cpp} (note 02)\\
\code{ISpectralFunction} & \code{benhar} & \code{nuclear::BenharSpectralFunction} & \file{physics/nuclear/SpectralFunctions.cpp} (note 02)\\
\code{IPhysicsModel} & \code{sigep} & \code{proton::SigEPModel} & \file{physics/proton/ProtonModels.cpp} (note 01)\\
\code{IPhysicsModel} & \code{deforest} & \code{proton::DeForestModel} & \file{physics/proton/ProtonModels.cpp} (note 02; flavour from \code{deForest\_flag})\\
\code{IPhysicsModel} & \code{deforest\_cc1} & \code{proton::DeForestCC1} & \file{physics/proton/ProtonModels.cpp} (note 02; $\sigma_{cc1}$ forced)\\
\code{IPhysicsModel} & \code{deforest\_cc2} & \code{proton::DeForestCC2} & \file{physics/proton/ProtonModels.cpp} (note 02; $\sigma_{cc2}$ forced)\\
\code{IPhysicsModel} & \code{deforest\_cc1\_onshell} & \code{proton::DeForestCC1OnShell} & \file{physics/proton/ProtonModels.cpp} (note 02)\\
\code{IPhysicsModel} & \code{nuc\_elastic} & \code{nuclear::NucElasticModel} & \file{physics/nuclear/NuclearElastic.cpp} (note 03)\\
\code{IRadiator} & \code{simc\_rad} & \code{SimcRadiator} & \file{radiation/SimcRadiator.cpp} (\secr{sec:radiation})\\
\code{ISpectrometer} & \code{hms} & \code{hms::HmsSpectrometer} & spectrometers/hms/ (HMS, beam right)\\
\code{ISpectrometer} & \code{shms} & \code{shms::ShmsSpectrometer} & spectrometers/shms/ (SHMS, beam left)\\
\code{ISpectrometer} & \code{sos} & \code{sos::SosSpectrometer} & spectrometers/sos/ (SOS, beam left)\\
\code{ISpectrometer} & \code{hrsr} & \code{hrs::HrsrSpectrometer} & spectrometers/hrs/ (HRS-R, beam right)\\
\code{ISpectrometer} & \code{hrsl} & \code{hrs::HrslSpectrometer} & spectrometers/hrs/ (HRS-L, beam left)\\
\code{ISpectrometer} & \code{calo} & \code{calo::CaloSpectrometer} & spectrometers/calo/ (calorimeter arms 7/8)\\
\code{IOutputSink} & \code{root\_tree} & \code{RootTreeSink} & \file{output/RootTreeSink.cpp} (the event tree \code{h10})\\
\code{IOutputSink} & \code{run\_summary} & \code{RunSummarySink} & \file{output/RunSummarySink.cpp} (summary JSON/text)\\
\code{IOutputSink} & \code{gen\_hist} & \code{GenHistSink} & \file{output/GenHistSink.cpp} (generated histograms)\\
\code{IOutputSink} & \code{fortran\_hist} & \code{driver::FortranHistSink} & \file{driver/ReportSink.cpp} (\texttt{.hist}, \texttt{.geni}, \texttt{.gen})\\
\end{simctable}

\subsection{How the pipeline selects the plug-ins}\lbl{sec:pipeline}
\code{PipelineConfig::Defaults} (driver/Pipeline.cpp) reproduces the choice of the Fortran from the input flags; the
first matching row of \tabr{tab:defaults} wins. The spectrometer plug-ins follow the arm codes \code{electron\_arm},
\code{hadron\_arm}; the random source is \code{mt64}, the target \code{hallc\_cryo}, the beam \code{standard}, the
radiator \code{simc\_rad}, and the sinks are \code{root\_tree} (when \code{Nntu} $>0$), \code{run\_summary},
\code{gen\_hist} and \code{fortran\_hist}.

\begin{simctable}{footnotesize}{L{0.36\linewidth}L{0.2\linewidth}L{0.17\linewidth}L{0.2\linewidth}}
\caption{Default reaction and physics model (\code{PipelineConfig::Defaults}; one rule per row; the reaction flags are
derived from the inputs by \code{DeriveReactionFlags}, core/RunSetup.cpp).}\lbl{tab:defaults}\\
\toprule\rowcolor{simchead}\simcth{condition} & \simcth{reaction} & \simcth{physics model} & \simcth{note}\\\midrule\endfirsthead
\toprule\rowcolor{simchead}\simcth{condition} & \simcth{reaction} & \simcth{physics model} & \simcth{note}\\\midrule\endhead
\bottomrule\endlastfoot
\code{doing\_phsp} = 1 & \code{phase\_space} & \code{sigep} (not used, $\sigma=1$) & --\\
no reaction flag, nint($A$) = 1 & \code{hydrogen\_elastic} & \code{sigep} & 01\\
no reaction flag, nint($A$) = 2 & \code{deuterium\_eep} & \code{deforest} & 02\\
no reaction flag, nint($A$) $\ge3$ & \code{nucleus\_eep} & \code{deforest} & 02\\
\code{doing\_nuc\_elast} = 1 & \code{nuclear\_elastic} & \code{nuc\_elastic} & 03\\
\code{doing\_pion} = 1 & \code{pion\_production} & \code{pion} & 04\\
\code{doing\_kaon} = 1 & \code{kaon\_production} & \code{kaon} & 05\\
\code{doing\_delta} = 1 & \code{delta\_production} & \code{delta} & 06\\
\code{doing\_semi} = 1 & \code{sidis\_production} & \code{sidis} & 07\\
\code{doing\_rho} = 1 & \code{rho\_production} & \code{rho} & 08\\
\end{simctable}

The optional \code{"pipeline"} object of a JSON input (or the file given with \code{-{}-pipeline}) overrides the
defaults, one stage per key (\tabr{tab:jsonkeys}). A stage is given as a name or as an object
\code{\{"type": name, <configuration keys>\}}; \code{simc++ -{}-print-pipeline} prints the resolved choice.

\begin{simctable}{footnotesize}{L{0.27\linewidth}L{0.3\linewidth}L{0.37\linewidth}}
\caption{Keys of the JSON \code{pipeline} object (\code{PipelineConfig::apply}; one key per row).}\lbl{tab:jsonkeys}\\
\toprule\rowcolor{simchead}\simcth{key} & \simcth{selects} & \simcth{example}\\\midrule\endfirsthead
\toprule\rowcolor{simchead}\simcth{key} & \simcth{selects} & \simcth{example}\\\midrule\endhead
\bottomrule\endlastfoot
\code{rng} & the random source & \code{"rng": \{"type": "ranlux"\}}\\
\code{target} & the target model & \code{"target": "hallc\_cryo"}\\
\code{beam} & the beam & \code{"beam": "standard"}\\
\code{reaction} & the reaction & \code{"reaction": "nucleus\_eep"}\\
\code{reaction.spectral\_function} & the spectral function of D/A(e,e\pr p) & \code{\{"type": "benhar", "file": "benharsf\_12.dat"\}}\\
\code{physics} & the cross-section model & \code{"physics": "deforest\_cc2"}\\
\code{radiator} & the radiative corrections & \code{"radiator": "simc\_rad"}\\
\code{spectrometers.hadron} & the hadron-arm model & \code{"hadron": "hms"}\\
\code{spectrometers.electron} & the electron-arm model & \code{"electron": \{"type": "shms"\}}\\
\code{sinks} & the output list (replaces the default) & \code{["root\_tree", "run\_summary"]}\\
\code{threads} & the number of event-loop threads & \code{4}\\
\end{simctable}

\subsection{What a reaction and a physics model implement}\lbl{sec:interfaces}
A reaction plug-in implements \code{IReaction} (\tabr{tab:ireaction}); the (e,e\pr p) family shares the base class
\code{reactions::EepReactionBase} (reactions/EepReactionBase.cpp), whose \code{completeEv}, \code{completeRecon},
\code{completeMain} and \code{genVolume} contain the branches of all the (e,e\pr p)-type reactions, and the derived
plug-ins implement \code{generate} and \code{spectralWeight}. A physics model implements \code{IPhysicsModel}
(\tabr{tab:imodel}): it receives the completed vertex event and returns the cross section in the units the reaction
expects. Replacing a model therefore needs no change of the reaction, as long as the new model keeps those units and
the kinematic conventions of this note.

\begin{simctable}{footnotesize}{L{0.18\linewidth}L{0.27\linewidth}L{0.55\linewidth}}
\caption{The methods of \code{IReaction} (include/simc/stages/IReaction.h; one method per row).}\lbl{tab:ireaction}\\
\toprule\rowcolor{simchead}\simcth{method} & \simcth{Fortran origin} & \simcth{what it must do}\\\midrule\endfirsthead
\toprule\rowcolor{simchead}\simcth{method} & \simcth{Fortran origin} & \simcth{what it must do}\\\midrule\endhead
\bottomrule\endlastfoot
\code{generate} & \file{event.f:generate} & draw the reaction's variables after the beam part, call \code{completeEv}, then radiate through the radiator (or copy the vertex event to the original event)\\
\code{completeEv} & \file{event.f:complete\_ev} & solve the kinematics at the vertex, compute the jacobian, the target traversals of the outgoing particles and the per-event radiative constants\\
\code{completeRecon} & \file{event.f:complete\_recon\_ev} & complete the reconstructed event from the spectrometer reconstruction\\
\code{completeMain} & \file{event.f:complete\_main} & set the spectral-function weight, call the physics model and assemble the event weight\\
\code{genVolume} & \file{simc.f:383}-401 & return the generation volume\\
\code{prepareCentral} & \file{simc.f:calculate\_central} & optional: set the central-event quantities the model reads\\
\code{beginTry} & -- & optional: try counter for messages\\
\code{finalTargetState} & -- & optional: target values changed during the run, for the report\\
\end{simctable}

\begin{simctable}{footnotesize}{L{0.2\linewidth}L{0.74\linewidth}}
\caption{The methods of \code{IPhysicsModel} (include/simc/stages/IPhysicsModel.h; one method per row).}\lbl{tab:imodel}\\
\toprule\rowcolor{simchead}\simcth{method} & \simcth{role}\\\midrule\endfirsthead
\toprule\rowcolor{simchead}\simcth{method} & \simcth{role}\\\midrule\endhead
\bottomrule\endlastfoot
\code{sigcc} & cross section at the vertex kinematics; enters the weight\\
\code{sigccRecon} & cross section at the reconstructed kinematics (diagnostic; 1 where the Fortran does not compute it)\\
\code{sigccMain} & vertex cross section for models that also write fields of the event (meson models)\\
\code{beginEvent} & optional: event counter for warning messages\\
\end{simctable}

% =====================================================================================================================
\section{The event chain in SIMC++}\lbl{sec:chain}

\subsection{Event records}\lbl{sec:records}
\begin{simctable}{footnotesize}{L{0.18\linewidth}L{0.2\linewidth}L{0.56\linewidth}}
\caption{The records of one event (struct \code{Event} / \code{MainEvent}; one record per row).}\lbl{tab:records}\\
\toprule\rowcolor{simchead}\simcth{record} & \simcth{symbol} & \simcth{content}\\\midrule\endfirsthead
\toprule\rowcolor{simchead}\simcth{record} & \simcth{symbol} & \simcth{content}\\\midrule\endhead
\bottomrule\endlastfoot
\code{vertex} & no mark ($E'$, $\nu$, \ldots) & kinematics at the scattering vertex, used for the cross section and the spectral function: after the photon of tail 1, before those of tails 2 and 3\\
\code{orig} & superscript ``orig'' & the particles leaving the vertex: the vertex event with the photons of tails 2 and 3 removed and with the beam energy before tail 1; equal to \code{vertex} without radiation\\
\code{main.SP} & superscript ``SP'' & the particles at the spectrometer entrance: after energy loss, Coulomb deceleration and multiple scattering in the target and windows\\
\code{main.FP} & -- & the focal-plane track of the spectrometer model\\
\code{main.RECON} & -- & the target quantities reconstructed by the spectrometer model ($=$ SP when the arm has no model)\\
\code{recon} & superscript ``rec'' & the reconstructed event: momenta and angles from the reconstruction, corrected for the most probable energy loss, completed by \code{completeRecon}\\
\end{simctable}

\subsection{One try}\lbl{sec:onetry}
\code{SimcRun::step} (driver/SimcRun.cpp) runs one try; \tabr{tab:chain} gives the order of the calls. A failed step
ends the try, which still counts as tried. The order of the calls and of every random-number draw is the order of the
Fortran, which makes an exact replay possible with the RANLUX source and \code{-{}-fortran-compat}.

\begin{center}
\begin{tikzpicture}[node distance=4.2mm and 3mm, every node/.style={simc/stage, text width=29mm, font=\sffamily\scriptsize}]
\node (g) {\textbf{IBeam::generate}\\beam, vertex position, $E_\mathrm{in}$};
\node[right=of g] (c) {\textbf{IReaction::generate}\\generated variables, \code{completeEv}};
\node[right=of c] (r) {\textbf{IRadiator::radiate}\\tails 1--3, \code{orig}};
\node[below=of r] (m) {\textbf{MonteCarlo::run}\\target effects, ISpectrometer, reconstruction};
\node[left=of m] (cr) {\textbf{IReaction::\\completeRecon}\\reconstructed kinematics};
\node[left=of cr] (cm) {\textbf{IReaction::\\completeMain}\\$S$, IPhysicsModel, weight};
\node[below=of cm, simc/ok] (p) {\textbf{cuts, sums, IOutputSink::fill}};
\draw[simc/flow] (g) -- (c); \draw[simc/flow] (c) -- (r); \draw[simc/flow] (r) -- (m);
\draw[simc/flow] (m) -- (cr); \draw[simc/flow] (cr) -- (cm); \draw[simc/flow] (cm) -- (p);
\end{tikzpicture}
\end{center}

\begin{simctable}{footnotesize}{L{0.05\linewidth}L{0.3\linewidth}L{0.38\linewidth}L{0.21\linewidth}}
\caption{The steps of one try (one step per row).}\lbl{tab:chain}\\
\toprule\rowcolor{simchead}\simcth{\#} & \simcth{C++ call} & \simcth{what happens} & \simcth{Fortran origin}\\\midrule\endfirsthead
\toprule\rowcolor{simchead}\simcth{\#} & \simcth{C++ call} & \simcth{what happens} & \simcth{Fortran origin}\\\midrule\endhead
\bottomrule\endlastfoot
1 & \code{IBeam::generate} & beam spot and raster, vertex $z$ (\code{ITarget::generateZ}), beam energy loss, beam energy at the vertex (\eqr{eq:ein}) & \file{event.f:147}-212\\
2 & \code{IReaction::generate} & the reaction's generated variables (\secr{sec:generation} and the reaction note) & \file{event.f:217}-373\\
3 & \code{IReaction::completeEv} & vertex kinematics, jacobian, traversals of the outgoing particles, \code{IRadiator::initEvent} & \file{event.f:432}-1054\\
4 & \code{IRadiator::radiate} & photons of the tails and the \code{orig} event; a tail-1 photon calls \code{completeEv} again & \file{event.f:398}-414, \file{radc.f:120}-528\\
5 & \code{MonteCarlo::run} & spectrometer-entrance quantities, \code{ISpectrometer::transport} of each arm, reconstruction, most-probable energy-loss correction & \file{simc.f:1320}-1862\\
6 & \code{IReaction::completeRecon} & reconstructed kinematics & \file{event.f:1058}-1362\\
7 & \code{IReaction::completeMain} & spectral-function weight, \code{IPhysicsModel::sigcc}, Coulomb factor, event weight & \file{event.f:1366}-1572\\
8 & \code{PassCuts}, hard cuts & acceptance cuts on the reconstructed event (\secr{sec:glossary}) & \file{simc.f:236}-253\\
9 & \code{IOutputSink::fill} & every try is offered to the sinks; the tree keeps the successes & \file{simc.f:255}-346\\
\end{simctable}

% =====================================================================================================================
\section{Notation and glossary}\lbl{sec:notation}

\subsection{Symbols}
The symbols of \tabr{tab:symbols} are used with the same meaning in all notes. Vertex quantities carry no mark;
reconstructed ones the superscript ``rec''.

\begin{simctable}{footnotesize}{L{0.17\linewidth}L{0.5\linewidth}L{0.09\linewidth}L{0.18\linewidth}}
\caption{Symbols (one symbol per row; the last column is the name in the code).}\lbl{tab:symbols}\\
\toprule\rowcolor{simchead}\simcth{symbol} & \simcth{meaning} & \simcth{unit} & \simcth{code}\\\midrule\endfirsthead
\toprule\rowcolor{simchead}\simcth{symbol} & \simcth{meaning} & \simcth{unit} & \simcth{code}\\\midrule\endhead
\bottomrule\endlastfoot
$E_b$ & nominal beam energy (input) & MeV & \code{Ebeam}\\
$\Delta E_b$ & full width of the flat beam-energy spread & MeV & \code{dEbeam}\\
$E_\mathrm{in}$ & beam energy at the vertex & MeV & \code{Ein}\\
$E_b^\mathrm{ave}$ & beam energy at the vertex for the reconstruction, $E_b+V_C-\Delta E_1^\mathrm{mp}$ & MeV & \code{Ebeam\_vertex\_ave}\\
$E'$, $P'_e$ & energy and momentum of the scattered electron ($P'_e=E'$) & MeV & \code{e.E}, \code{e.P}\\
$\theta_e$, $\varphi_e$ & polar and azimuthal angle of the scattered electron & rad & \code{e.theta}, \code{e.phi}\\
$E_h$, $P_h$, $M_h$ & energy, momentum and mass of the detected hadron & MeV & \code{p.E}, \code{p.P}, \code{Mh}\\
$\theta_h$, $\varphi_h$ & polar and azimuthal angle of the hadron & rad & \code{p.theta}, \code{p.phi}\\
$P^0_e$, $\theta^0_e$ & central momentum and angle of the electron arm & MeV, rad & \code{spec.e.P}, \code{spec.e.theta}\\
$P^0_h$, $\theta^0_h$ & central momentum and angle of the hadron arm & MeV, rad & \code{spec.p.P}, \code{spec.p.theta}\\
$\varphi_0$ & central azimuth of an arm ($3\pi/2$ beam right, $\pi/2$ beam left) & rad & \code{spec.e/p.phi}\\
$x'$, $y'$ & slopes of a track relative to the central ray (\secr{sec:specvars}) & rad & \code{xptar}, \code{yptar}\\
$\delta$ & relative momentum, $100(P-P^0)/P^0$ & \% & \code{delta}\\
$M_A$ & target mass & MeV & \code{targ.M}\\
$M_N$ & mass of the struck nucleon & MeV & \code{targ.Mtar\_struck}\\
$M_\mathrm{rec}$ & mass of the residual system (A$-$1) & MeV & \code{targ.Mrec}\\
$\nu$, $Q^2$, $\vec q$ & energy transfer, four-momentum transfer squared ($Q^2>0$), three-momentum transfer & MeV, MeV$^2$ & \code{nu}, \code{Q2}, \code{q}\\
$x$, $W$, $\varepsilon$ & Bjorken $x$, invariant mass of the photon--nucleon system, virtual-photon polarisation & --, MeV, -- & \code{xbj}, \code{W}, \code{epsilon}\\
$\vec P_m$, $P_m$ & missing momentum $\vec p_h-\vec q$ and its magnitude & MeV/$c$ & \code{Pmx/y/z}, \code{Pm}\\
$E_m$ & missing energy (definition per reaction) & MeV & \code{Em}\\
$T_\mathrm{rec}$ & kinetic energy of the residual system & MeV & \code{Trec}\\
$V_C$ & Coulomb energy of the electron at the nucleus & MeV & \code{targ.Coulomb}\\
$\Delta E_i$ & energy loss of particle $i$ (1 beam, 2 electron, 3 hadron); mp: most probable & MeV & \code{Eloss(i)}\\
$t_i$ & material crossed by particle $i$, in radiation lengths & -- & \code{teff(i)}\\
$E'_\mathrm{min}$, $E'_\mathrm{max}$ & limits of the scattered-electron energy at the vertex (edges, \secr{sec:edges}) & MeV & \code{edge.e.E}\\
$E_{h,\mathrm{min}}$, $E_{h,\mathrm{max}}$ & limits of the hadron energy at the vertex (edges) & MeV & \code{edge.p.E}\\
$\Sigma_\mathrm{min}$, $\Sigma_\mathrm{max}$ & limits of the summed generated energy & MeV & \code{gen.sumEgen}\\
$d_E$ & energy margin added to the edges (input) & MeV & \code{dE\_edge\_test}\\
$E_{\gamma i}$ & energy of the photon radiated by tail $i$ & MeV & \code{Egamma\_used(i)}\\
$E_\gamma^\mathrm{tot,max}$ & largest total photon energy & MeV & \code{Egamma\_tot\_max}\\
$V_\mathrm{gen}$ & generation volume & sr$^k$ MeV$^l$ & \code{genvol}\\
$\mathcal L$ & luminosity & $\mu$b$^{-1}$ & \code{luminosity}\\
$J$ & jacobian of the event & -- & \code{main.jacobian}\\
$w_\mathrm{gen}$ & generation and radiative weight of the event & -- & \code{main.gen\_weight}\\
$S$ & spectral-function weight & MeV$^{-3}$ or MeV$^{-4}$ & \code{main.SF\_weight}\\
$\sigma$ & cross section at the vertex & per reaction & \code{main.sigcc}\\
$w$ & event weight & per reaction & \code{main.weight}\\
\end{simctable}

\subsection{Glossary}\lbl{sec:glossary}
\begin{simctable}{footnotesize}{L{0.2\linewidth}L{0.76\linewidth}}
\caption{Terms of SIMC used in the notes (one term per row).}\lbl{tab:glossary}\\
\toprule\rowcolor{simchead}\simcth{term} & \simcth{definition}\\\midrule\endfirsthead
\toprule\rowcolor{simchead}\simcth{term} & \simcth{definition}\\\midrule\endhead
\bottomrule\endlastfoot
try & one pass through the event chain (\tabr{tab:chain}); counted in $N_\mathrm{tried}$ (\code{ntried}) whether or not it succeeds\\
success & a try that passes every step: generation, radiation, both spectrometers, reconstruction, a positive spectral-function weight and, with hard cuts, the cuts. Successes are counted in $N_\mathrm{event}$ (\code{nevent}), written to the tree and fill the reconstructed histograms. A run stops after \code{ngen} successes (or $|\code{ngen}|$ tries if \code{ngen} $<0$)\\
contributing event & a success that also passes \code{pass\_cuts}; only contributing events enter \code{wtcontr} and the resolution sums\\
slope & $x'=\mathrm{d}x/\mathrm{d}z$ (\code{xptar}, vertical, $x$ down, the dispersive direction) and $y'=\mathrm{d}y/\mathrm{d}z$ (\code{yptar}, horizontal) of a track relative to the central ray of its arm, at the target\\
delta & $\delta = 100(P-P^0)/P^0$, in \%\\
acceptance & the input limits \code{SPedge} of $\delta$, $y'$, $x'$ of each arm (slopes entered in mrad)\\
slop & a margin added to a limit: the spectrometer slops (\code{slop\_param\_*}, the resolution of each arm's reconstruction, hard-coded per spectrometer), and the total slop of the missing energy\\
edge & a limit at an intermediate stage, derived from the acceptance by adding slops, energy losses, Coulomb shifts and multiple-scattering angles (\secr{sec:edges})\\
generation limits & the ranges in which the generated variables are drawn (\code{gen}); for the slopes they equal the edges\\
pass\_cuts & the test that the reconstructed $\delta$, $y'$, $x'$ of both arms lie inside the acceptance shrunk by the spectrometer slops\\
hard cuts & input \code{hard\_cuts} = 1: events failing \code{pass\_cuts}, and (e,e\pr p)-type events with $E_m^\mathrm{rec}$ above \code{cuts\%Em\%max}, are not successes\\
generation volume & $V_\mathrm{gen}$ (\code{genvol}), the product of the full generation ranges (\eqr{eq:genvol})\\
normfac & $\mathcal L V_\mathrm{gen}N_\mathrm{event}/N_\mathrm{tried}$ (\eqr{eq:normfac}); weight a tree histogram with $w\cdot\mathrm{normfac}/N_\mathrm{event}$ to obtain counts\\
wtcontr & the expected number of counts of the simulated charge inside the cuts, printed in the \texttt{.hist} (\eqr{eq:normfac})\\
Weight & the tree column holding the event weight $w$ (\eqr{eq:weight})\\
tail & one of the three radiating legs: 1 incoming electron, 2 scattered electron, 3 outgoing hadron\\
central event & the event at the central momenta and angles of both arms, used for the \texttt{.hist} line \code{CENTRAL.sigcc}\\
default mode & SIMC++ with the clear Fortran bugs corrected\\
compat mode & SIMC++ with \code{-{}-fortran-compat}: the exact Fortran behaviour\\
\end{simctable}

\subsection{Constants}
\begin{simctable}{footnotesize}{L{0.2\linewidth}L{0.34\linewidth}L{0.4\linewidth}}
\caption{Constants (core/SimcConstants.h; from \fsrc{constants.inc:11-47}; one constant per row).}\lbl{tab:constants}\\
\toprule\rowcolor{simchead}\simcth{symbol} & \simcth{value} & \simcth{code}\\\midrule\endfirsthead
\toprule\rowcolor{simchead}\simcth{symbol} & \simcth{value} & \simcth{code}\\\midrule\endhead
\bottomrule\endlastfoot
$m_e$ & 0.51099906 MeV & \code{Me}\\
$M_p$ & 938.27231 MeV & \code{Mp}\\
$M_n$ & 939.56563 MeV & \code{Mn}\\
$M_d$ & 1875.613 MeV & \code{Md}\\
$u$ & 931.49432 MeV (atomic mass unit) & \code{amu}\\
$\hbar c$ & 197.327053 MeV\,fm & \code{hbarc}\\
$\alpha$ & 1/137.0359895 & \code{alpha}\\
$\gamma_E$ & 0.577215665 & \code{euler}\\
\end{simctable}
The bremsstrahlung routines work in GeV with their own copies of $m_e$, $M_p$ and $\alpha$ (\secr{sec:brem}).

% =====================================================================================================================
\section{Frames and sign conventions}\lbl{sec:frames}

\subsection{The laboratory frame}\lbl{sec:labframe}
The SIMC laboratory frame has $z$ along the beam (downstream), $x$ pointing down and $y$ pointing to the left looking
downstream. A direction is given by its polar angle $\theta$ to the beam and its azimuth $\varphi$ counted from $x$:
\begin{feq}{reactions/EepReactionBase.cpp: EepReactionBase::completeEv}{event.f:504-510}
\hat u = (\sin\theta\cos\varphi,\ \sin\theta\sin\varphi,\ \cos\theta).
\lbl{eq:unitvector}
\end{feq}
The arms on the beam right (HMS, HRS-R, calorimeter code 7) have $\varphi_0 = 3\pi/2$, those on the beam left (SOS,
HRS-L, SHMS, calorimeter code 8) $\varphi_0 = \pi/2$; $\theta_0$ is the input angle made positive
(core/SpectrometerId.cpp; \fsrc{dbase.f:243-268}). A code comment of the Fortran describes the beam-spot coordinates
\code{main\%target\%x,y} with $y$ up and $x$ to the right \fsrc{event.f:151-154}, which is not the orientation of
the frame above; the spot is symmetric in the default configuration, and the tree column \code{fry} is minus the
vertical raster coordinate.

\subsection{Spectrometer variables}\lbl{sec:specvars}
Each arm describes a track by $\delta$ and the slopes $x'$, $y'$ at the target (\tabr{tab:glossary}). For an electron
$P = E$ (the electron mass is neglected in the kinematics). The physics angles follow from the slopes with
$r=\sqrt{1+x'^2+y'^2}$:
\begin{feqa}{setup/KinematicsUtils.cpp: PhysicsAngles}{event.f:1594-1614}
\theta &= \arccos\frac{\cos\theta_0 - y'\sin\theta_0\sin\varphi_0}{r},\nonumber\\
\varphi &= \arctan\frac{y'\cos\theta_0 + \sin\theta_0\sin\varphi_0}{x'}
 \quad(+\pi \text{ if } \varphi\le 0;\ +\pi \text{ if } \sin\varphi_0<0;\ \varphi=\varphi_0 \text{ if } x'=0).
\lbl{eq:physang}
\end{feqa}
The formula assumes $\varphi_0 = \pi/2$ or $3\pi/2$, which all arms satisfy. The inverse mapping, used when a
direction is computed instead of generated, is
\begin{feqa}{setup/KinematicsUtils.cpp: SpectrometerAngles}{event.f:1636-1648}
\cos\Delta\theta &= \hat u\cdot\hat u_0,\qquad x' = \frac{u_x}{\cos\Delta\theta},\qquad
|y'| = \sqrt{\frac{1}{\cos^2\Delta\theta}-1-x'^2},
\lbl{eq:specang}
\end{feqa}
with $y'<0$ when $u_y/\cos\Delta\theta < u_{0,y}$. A rounding-negative argument of the square root gives NaN in the
Fortran; the default mode of SIMC++ sets $y'=0$ instead.

\paragraph{Coordinates at the target.} Before a spectrometer model is called, the vertex is written in the arm's
frame ($x$ down, $y$ horizontal, $z$ along the central ray) and drifted back to the plane $z=0$:
\begin{feqa}{driver/MonteCarlo.cpp: MonteCarlo::run}{simc.f:1421-1453, 1690-1720}
x_\mathrm{arm} &= -y_t - x_\mathrm{off},\qquad
y_\mathrm{arm} = -x_t\cos\theta_0 - z_t\sin\theta_0\sin\varphi_0 - y_\mathrm{off},\nonumber\\
z_\mathrm{arm} &= z_t\cos\theta_0 + x_t\sin\theta_0\sin\varphi_0 - z_\mathrm{off},\nonumber\\
x_\mathrm{arm} &\leftarrow x_\mathrm{arm} - z_\mathrm{arm}\,x',\qquad
y_\mathrm{arm} \leftarrow y_\mathrm{arm} - z_\mathrm{arm}\,y',\qquad z_\mathrm{arm}\leftarrow 0,
\lbl{eq:transport}
\end{feqa}
with $(x_t,y_t,z_t)$ the vertex position and the offsets those of the inputs \code{spec\%e(p)\%offset} (the slope
offsets are subtracted from the slopes).

\subsection{The angle \texorpdfstring{$\varphi_{pq}$}{phi\_pq}}\lbl{sec:phipq}
$\varphi_{pq}$ is defined in the Hall C analysis frame, $x_\mathrm{a}=-y$, $y_\mathrm{a}=x$, $z_\mathrm{a}=z$. With
$\hat q$ and $\hat p$ the directions of $\vec q$ and of the hadron in that frame, a frame $z'\parallel\hat q$,
$y'\parallel\hat q\times\hat z$, $x'=y'\times\hat q$ is built and
\begin{feq}{reactions/EepReactionBase.cpp: EepReactionBase::completeRecon}{event.f:1161-1188}
\cos\varphi_{pq}\propto\hat p\cdot\hat x',\qquad \sin\varphi_{pq}\propto\hat p\cdot\hat y',\qquad
\varphi_{pq}\in[0,2\pi),
\lbl{eq:phipq}
\end{feq}
(\code{atan2}$(\hat p\cdot\hat y',\hat p\cdot\hat x')$, shifted by $2\pi$ when negative); $\varphi_{pq}=0$ for a hadron in the scattering plane on
the downstream side of $\vec q$. The reconstructed value is computed for every reaction; the vertex value only by the
meson reactions.

% =====================================================================================================================
\section{Kinematic variables}\lbl{sec:kinvars}
The same formulas are applied at the vertex (\code{completeEv}) and to the reconstructed event (\code{completeRecon});
they differ in the beam energy and in the definitions noted below.

\paragraph{Beam energy.} At the vertex $E_\mathrm{in}$ is the generated beam energy after the energy loss and the
Coulomb acceleration of the incoming electron (\eqr{eq:ein}), lowered by the photon of tail 1 if there is one. The
reconstruction uses a fixed value that does not include the Coulomb shift:
\begin{feq}{reactions/EepReactionBase.cpp: EepReactionBase::completeRecon}{event.f:1089, 1098; init.f:281}
E_\mathrm{in}^\mathrm{rec} = E_b^\mathrm{ave} - V_C = E_b - \Delta E_1^\mathrm{mp}.
\lbl{eq:einrec}
\end{feq}

\paragraph{Electron side.}
\begin{feqa}{reactions/EepReactionBase.cpp: completeEv, completeRecon}{event.f:530-537, 1112-1117}
\nu &= E_\mathrm{in} - E',\qquad Q^2 = 2E_\mathrm{in}E'(1-\cos\theta_e),\qquad |\vec q| = \sqrt{Q^2+\nu^2},
\lbl{eq:nuq2}\\
\hat q &= \bigl(-E'\,u_{e,x},\ -E'\,u_{e,y},\ E_\mathrm{in}-E'\,u_{e,z}\bigr)/|\vec q|,
\lbl{eq:qhat}\\
x &= Q^2/(2M_p\nu).\lbl{eq:xbj}
\end{feqa}
The run stops if $\hat q$ is not a unit vector to 1\,\% \fsrc{event.f:538-539}. $x$ uses $M_p$ for every target.

\paragraph{Invariant mass and polarisation.}
\begin{feqa}{reactions/EepReactionBase.cpp: completeRecon; reactions/pion/PionProduction.cpp: completeEv}{event.f:724-726, 1126-1128, 1142}
W^2 &= M_N^2+2M_N\nu-Q^2\ \ (\text{vertex, meson reactions}),\qquad
(W^\mathrm{rec})^2 = M_p^2+2M_p\nu^\mathrm{rec}-(Q^\mathrm{rec})^2,\nonumber\\
W &= \operatorname{sign}(W^2)\sqrt{|W^2|},
\lbl{eq:w}\\
\varepsilon &= \Bigl[1+2\bigl(1+\nu^2/Q^2\bigr)\tan^2(\theta_e/2)\Bigr]^{-1}.\lbl{eq:epsilon}
\end{feqa}
The reconstructed $W$ always uses the proton mass, as the Hall C analysis does. The vertex values of $W$,
$\varepsilon$, $\theta_{pq}$, $\varphi_{pq}$ and $t$ are computed only by the meson reactions; for the (e,e\pr p)-type
reactions they stay 0 (so are the tree columns \code{Wi} and \code{epsiloni}).

\paragraph{Hadron side.} $\theta_{pq} = \arccos(\hat p\cdot\hat q)$ (argument limited to 1 in the reconstruction).
The missing-momentum vector and the missing energy of the whole target are
\begin{feq}{reactions/EepReactionBase.cpp: completeEv, completeRecon}{event.f:914-918, 1312-1315}
\vec P_m = P_h\,\hat p - |\vec q|\,\hat q,\qquad P_\mathrm{miss}=|\vec P_m|,\qquad E_\mathrm{miss} = \nu + M_A - E_h,
\lbl{eq:pm}
\end{feq}
$\vec P_m=\vec p_h-\vec q$ is the momentum of the struck nucleon before the collision in the impulse approximation;
the residual system recoils with $-\vec P_m$. Its components along $\hat q$, along the out-of-plane direction
$(-\hat q_y,\hat q_x,0)/\sqrt{\hat q_x^2+\hat q_y^2}$ and the remainder are
\begin{feq}{reactions/EepReactionBase.cpp: completeEv, completeRecon}{event.f:922-926, 1319-1323}
P_{m,\parallel} = \vec P_m\cdot\hat q,\quad
P_{m,\mathrm{oop}} = \frac{-P_{m,x}\hat q_y + P_{m,y}\hat q_x}{\sqrt{\hat q_x^2+\hat q_y^2}},\quad
P_{m,\perp} = \sqrt{\max\bigl(0,\ P_m^2-P_{m,\parallel}^2-P_{m,\mathrm{oop}}^2\bigr)}\ \ (\ge0).
\lbl{eq:pmcomp}
\end{feq}
Looking along $\vec q$, the out-of-plane direction points down; the sign of $P_{m,\perp}$ is not determined. $E_m$,
the scalar $P_m$ used by a reaction and $T_\mathrm{rec}$ are defined in each reaction note. Every reaction also
computes $z = E_h/\nu$ and $p_T^2 = P_h^2\sin^2\theta_{pq}$.

% =====================================================================================================================
\section{Beam, vertex and generation}\lbl{sec:generation}

\subsection{Beam and vertex}\lbl{sec:beam}
The beam spot is a Gaussian truncated at $3\sigma$ (widths \code{gen\%xwid}, \code{gen\%ywid}) plus the raster
($\xi_i$ uniform random numbers, $f_{1,2}$ = \code{targ\%fr1,fr2}, $L$ the target length):
\begin{feqa}{reactions/StandardBeam.cpp: StandardBeam::generate; material/HallCCryoTarget.cpp: generateZ}{event.f:156-187}
\text{pattern 1:}&\quad (x_r,y_r) = (f_1\cos\pi\xi_1,\ f_2\cos\pi\xi_2),\nonumber\\
\text{pattern 2:}&\quad (x_r,y_r)=\rho\,(\cos2\pi\xi_1,\ \sin2\pi\xi_1),\quad \rho=\sqrt{\xi_2}\,(f_2-f_1)+f_1,\nonumber\\
\text{pattern 3:}&\quad (x_r,y_r) = \bigl(f_1(2\xi_1-1),\ f_2(2\xi_2-1)\bigr),\nonumber\\
&\quad z_t = (0.5-\xi_3)\,L + z_\mathrm{off}.\lbl{eq:raster}
\end{feqa}
The beam energy at the vertex is
\begin{feq}{reactions/StandardBeam.cpp: StandardBeam::generate}{event.f:194-205}
E_\mathrm{in} = E_b + (\xi-0.5)\,\Delta E_b + V_C - \Delta E_1,
\lbl{eq:ein}
\end{feq}
with $\Delta E_b$ = \code{Ebeam}$\cdot$\code{dEbeam}/100, $\Delta E_1$ the sampled energy loss of the beam up to $z_t$
(0 without \code{using\_Eloss}) and $V_C$ the Coulomb energy (0 without \code{using\_Coulomb}). The random numbers for
the spread and the energy loss are drawn even when they are not used, so that the random stream does not depend on
these switches.

\subsection{Edges}\lbl{sec:edges}
The limits used for the generation are derived once at start-up from the acceptance of each arm, by widening it with
the spectrometer slops (only for an arm with a model), and propagating it to the vertex:
\begin{feqa}{setup/SetupLimits.cpp: BuildLimits}{init.f:209-232, 260-271}
E'_\mathrm{min} &= (1+\delta^\mathrm{acc}_\mathrm{min}/100)P_e^0 + V_C - d_E + \Delta E_2^\mathrm{min},\nonumber\\
E'_\mathrm{max} &= (1+\delta^\mathrm{acc}_\mathrm{max}/100)P_e^0 + V_C + d_E + \Delta E_2^\mathrm{max},\nonumber\\
E_{h,\mathrm{min/max}} &= \sqrt{p_\mathrm{min/max}^2+M_h^2}+\Delta E_3^\mathrm{min/max},\nonumber\\
p_\mathrm{min/max} &=(1+\delta^\mathrm{acc}_\mathrm{min/max}/100)P_h^0\mp d_E\ \ (p_\mathrm{min}\ge0.001),\nonumber\\
x'_\mathrm{min} &= x'^\mathrm{acc}_\mathrm{min}-\theta^\mathrm{ms}_\mathrm{max},\qquad
x'_\mathrm{max} = x'^\mathrm{acc}_\mathrm{max}+\theta^\mathrm{ms}_\mathrm{max}\quad(\text{same for }y'),
\lbl{eq:edges}
\end{feqa}
where the acceptance limits already include the slops, $\Delta E_i^\mathrm{min/max}$ are the extreme energy losses of
particle $i$ in the target (\code{extremeLosses}, \fsrc{target.f:310-544}) and $\theta^\mathrm{ms}_\mathrm{max}$ is
$3.5$ times the r.m.s.\ multiple-scattering angle for the thickest path. The generated slopes are drawn inside these
edges. The beam-energy limits are
\begin{feq}{setup/SetupLimits.cpp: BuildLimits}{init.f:281, 285-286}
E_b^\mathrm{max/min} = E_b\pm\Delta E_b/2 - \Delta E_1^\mathrm{min/max} + V_C,\qquad
E_b^\mathrm{ave} = E_b+V_C-\Delta E_1^\mathrm{mp},
\lbl{eq:ebeamlimits}
\end{feq}
and the limits on the vertex missing energy and momentum, on the summed generated energy $\Sigma$, on the generated
energies and on the photon energies follow from energy conservation per reaction (reaction notes, \S2).

\subsection{Flat generation and the generation volume}
Every generated variable is drawn flat in its range, $v = v_\mathrm{min}+\xi(v_\mathrm{max}-v_\mathrm{min})$. When the
range of an energy is narrowed event by event (by the summed-energy limits), the event receives the factor
\begin{feq}{reactions/NuclearEep.cpp: NuclearEep::generate}{event.f:293, 312}
w_\mathrm{gen} \leftarrow w_\mathrm{gen}\,\frac{E_\mathrm{hi}-E_\mathrm{lo}}{E^\mathrm{gen}_\mathrm{max}-E^\mathrm{gen}_\mathrm{min}},
\lbl{eq:genweight}
\end{feq}
so that the generation volume is the product of the full ranges:
\begin{feqa}{reactions/EepReactionBase.cpp: EepReactionBase::genVolume}{simc.f:383-401}
\Delta\Omega_e &= (y'_{e,\mathrm{max}}-y'_{e,\mathrm{min}})(x'_{e,\mathrm{max}}-x'_{e,\mathrm{min}}),\qquad
\Delta\Omega_h = (y'_{h,\mathrm{max}}-y'_{h,\mathrm{min}})(x'_{h,\mathrm{max}}-x'_{h,\mathrm{min}}),\nonumber\\
V_\mathrm{gen} &= \Delta\Omega_e\ \times\bigl[\Delta\Omega_h\,\Delta E'\bigr]\times\bigl[\Delta E_h\bigr],
\lbl{eq:genvol}
\end{feqa}
the first bracket for every reaction except H(e,e\pr p) and A(e,e\pr)A, the second for A(e,e\pr p) and SIDIS
($\Delta\Omega$ in sr, $\Delta E$ in MeV; $4\pi$ replaces $\Delta\Omega_h$ for the $\rho$).

\subsection{The solid-angle jacobian}\lbl{sec:jacobian}
The slopes are generated flat, but the cross sections are per unit solid angle, and
$\mathrm{d}\Omega = \mathrm{d}x'\,\mathrm{d}y'/r^3$. Each generated pair of slopes gives the factor
\begin{feq}{reactions/EepReactionBase.cpp: EepReactionBase::completeEv}{event.f:1015-1025}
J_\Omega = r^{-3},\qquad r = \sqrt{1+x'^2+y'^2} = 1/\cos(\theta-\theta_0),
\lbl{eq:jacomega}
\end{feq}
applied for the electron in every reaction and for the hadron in A(e,e\pr p) and the meson reactions (not for D(e,e\pr p),
see note 02). Reaction-specific jacobians multiply into the same $J$.

% =====================================================================================================================
\section{Event weight and normalisation}\lbl{sec:weight}

\subsection{The event weight}
\begin{feq}{reactions/EventChain.cpp: assembleWeight}{event.f:1561-1564}
w = S\cdot J\cdot w_\mathrm{gen}\cdot\sigma\cdot n_\mathrm{tgt}\ [\cdot P_\mathrm{surv}],
\lbl{eq:weight}
\end{feq}
with $S$ the spectral-function weight (1 for reactions without one), $J$ the jacobian, $w_\mathrm{gen}$ the generation
weight (\eqr{eq:genweight}) times the radiative weight (\eqr{eq:radweight}), $\sigma$ the cross section at the vertex
times the Coulomb factor (\eqr{eq:coulombfocus}), $n_\mathrm{tgt}$ the number of nucleons that can take part (meson
reactions; 1 otherwise) and $P_\mathrm{surv}$ the survival probability of a kaon whose decay is not simulated. An
event with $S\le0$ is not a success. The tree column \code{Weight} is $w$, the column \code{sigcc} is $\sigma$.

\subsection{Luminosity and normalisation}\lbl{sec:norm}
\begin{feqa}{setup/SetupBuilder.cpp: BuildRunSetup}{simc.f:101-108}
f_\mathrm{tgt} &= \frac{A_\mathrm{amu}}{3.75914\times10^{6}\cdot(a/100)}\cdot\frac{|\cos\theta_\mathrm{tgt}|}{t\,[\mathrm{mg/cm^2}]},
\qquad \mathcal{L} = \frac{Q\,[\mathrm{mC}]}{f_\mathrm{tgt}}\quad[\mu\mathrm{b}^{-1}],
\lbl{eq:lumi}
\end{feqa}
with $A_\mathrm{amu}$ = \code{targ\%mass\_amu}, $a$ = \code{targ\%abundancy} (\%), $t$ = \code{targ\%thick},
$\theta_\mathrm{tgt}$ = \code{targ\%angle} and $Q$ = \code{EXPER\%charge}. The constant is (electrons per mC)$\times
N_A\times$(mg$\to$g)$\times$(cm$^2\to\mu$b); the CODATA 2018 value of this product is $3.758724\times10^6$, so the
hard-coded value is larger by $1.1\times10^{-4}$. The quotient is evaluated in quadruple precision, as in the Fortran
build. The target length is $L=t/\rho$. After the loop
\begin{feq}{driver/RunReport.cpp: Normalise}{simc.f:375, 403-405, 930}
\mathrm{normfac} = \frac{\mathcal{L}\,V_\mathrm{gen}\,N_\mathrm{event}}{N_\mathrm{tried}}\ (1\text{ for phase space}),\qquad
\mathrm{wtcontr} = \frac{\mathcal{L}\,V_\mathrm{gen}}{N_\mathrm{tried}}\sum_{\text{contributing}} w_i.
\lbl{eq:normfac}
\end{feq}
$\sum w_i/N_\mathrm{tried}$ is the Monte Carlo estimate of the mean of $w$ over the generation volume, so
\code{wtcontr} is the number of counts of the simulated experiment inside the cuts.

% =====================================================================================================================
\section{Target effects on the scattered particles}\lbl{sec:spquantities}
The true (\code{orig}) particles become the quantities at the spectrometer entrance:
\begin{feqa}{driver/MonteCarlo.cpp: MonteCarlo::run}{simc.f:1389-1409, 1665-1678}
\delta_h^\mathrm{SP} &= 100\,\frac{\sqrt{|(E_h-\Delta E_3)^2-M_h^2|}-P_h^0}{P_h^0}\ \ (\delta_h\text{ without energy loss}),\nonumber\\
\delta_e^\mathrm{SP} & = 100\,\frac{E'-\Delta E_2-V_C-P_e^0}{P_e^0},\nonumber\\
y'^\mathrm{SP} &= y' + \theta^\mathrm{ms}_1 + \theta^\mathrm{ms,beam}_1,\qquad x'^\mathrm{SP} = x' + \theta^\mathrm{ms}_2 + \theta^\mathrm{ms,beam}_2\cos\theta_0,
\lbl{eq:sp}
\end{feqa}
with $\Delta E_{2,3}$ sampled in \code{completeEv} and the multiple-scattering angles of \eqr{eq:musc} (only with
\code{mc\_smear}; the beam angles, drawn with $t_1$, are added to both arms). The electron loses $V_C$ when it leaves
the nucleus; the hadron is not affected by the Coulomb field. After the spectrometer model the momenta are rebuilt from
the reconstructed $\delta$ and, with \code{correct\_Eloss}, the most probable losses at the target centre are added back:
\begin{feqa}{driver/MonteCarlo.cpp: MonteCarlo::run}{simc.f:1633-1655, 1829-1853}
P^\mathrm{rec} &= P^0(1+\delta^\mathrm{rec}/100),\nonumber\\
E_h^\mathrm{rec} &= \max\Bigl(\sqrt{(P_h^\mathrm{rec})^2+M_h^2}+\Delta E_3^\mathrm{mp},\ \sqrt{M_h^2+10^{-6}}\Bigr),\nonumber\\
E'^\mathrm{rec} &= E'^\mathrm{rec}_0 + \Delta E_2^\mathrm{mp}.
\lbl{eq:recon}
\end{feqa}
The reconstructed $\delta$ itself is not corrected, and the electron energy is not corrected for $V_C$.

% =====================================================================================================================
\section{Energy loss, radiation length and multiple scattering}\lbl{sec:eloss}

\subsection{Paths through the target and the windows}
\code{HallCCryoTarget::tripThruTarget} returns the energy loss and the thickness in radiation lengths seen by the beam
(1), the scattered electron (2) or the hadron (3) for a vertex at $z$ from the target centre. A target is ``liquid'' when
$Z<2.4$. The beam crosses $(L/2+z)/|\cos\theta_\mathrm{tgt}|$ of target and, for a liquid, the cell entrance window
(\tabr{tab:cells}); a scattered particle crosses $(L/2-z)/|\cos(\theta\pm\theta_\mathrm{tgt})|$ of target (or the shorter
path to the cell wall of a liquid cell), the cell wall, the exit windows and the air of its arm (\tabr{tab:windows}).

\begin{simctable}{footnotesize}{L{0.2\linewidth}L{0.22\linewidth}L{0.52\linewidth}}
\caption{Cryogenic cells (\code{targ\%can}; \fsrc{target.f:34-44, 106-168}; one cell per row).}\lbl{tab:cells}\\
\toprule\rowcolor{simchead}\simcth{\code{targ\%can}} & \simcth{entrance window} & \simcth{side wall for the scattered particles}\\\midrule\endfirsthead
\bottomrule\endlastfoot
1 (beer can) & 2.8 mil Al & radius 1.325\,in, 5 mil Al; the shorter of forward and side path\\
2 (pudding can) & 5 mil Al & spherical end, 5 mil Al, intersection computed per event\\
3 (2017 10-cm cell) & 0.013 cm Al & cylinder with round end cap, inner radius 1.315\,in, wall 7.1 mil\\
\end{simctable}

\begin{simctable}{footnotesize}{L{0.17\linewidth}L{0.15\linewidth}L{0.12\linewidth}L{0.12\linewidth}L{0.12\linewidth}L{0.22\linewidth}}
\caption{Windows and air of each arm (\fsrc{target.f:73-103, 189-223}; 1 mil = $2.54\times10^{-3}$ cm; one arm per row).}\lbl{tab:windows}\\
\toprule\rowcolor{simchead}\simcth{arm} & \simcth{Al} & \simcth{air} & \simcth{kevlar} & \simcth{mylar} & \simcth{path angle}\\\midrule\endfirsthead
\toprule\rowcolor{simchead}\simcth{arm} & \simcth{Al} & \simcth{air} & \simcth{kevlar} & \simcth{mylar} & \simcth{path angle}\\\midrule\endhead
\bottomrule\endlastfoot
HMS (1) & 16 mil & 15 cm & 15 mil & 5 mil & $\theta+\theta_\mathrm{tgt}$\\
SOS (2) & 8 mil & 15 cm & 5 mil & 3 mil & $\theta-\theta_\mathrm{tgt}$\\
HRS-R (3) & 13 mil & 15 cm & -- & 10 mil & $\theta-\theta_\mathrm{tgt}$\\
HRS-L (4) & 13 mil & 15 cm & -- & 10 mil & $\theta-\theta_\mathrm{tgt}$\\
SHMS (5) & 30 mil & 57.27 cm & -- & -- & $\theta-\theta_\mathrm{tgt}$\\
\end{simctable}

\begin{simctable}{footnotesize}{L{0.15\linewidth}L{0.15\linewidth}L{0.1\linewidth}L{0.1\linewidth}L{0.15\linewidth}}
\caption{Window materials (\fsrc{target.inc:5-29}; one material per row).}\lbl{tab:materials}\\
\toprule\rowcolor{simchead}\simcth{material} & \simcth{$\rho$ [g/cm$^3$]} & \simcth{$Z$} & \simcth{$A$} & \simcth{$X_0$ [g/cm$^2$]}\\\midrule\endfirsthead
\bottomrule\endlastfoot
aluminium & 2.70 & 13 & 26.98 & 24.01\\
mylar & 1.39 & 4.545 & 8.735 & 39.95\\
kevlar & 0.74 & 2.67 & 4.67 & 55.2\\
air & 0.00121 & 7.2 & 14.4 & 36.66\\
\end{simctable}
The calorimeter arms 7/8 have no window description of their own (they keep the values of the previous arm).

\subsection{Radiation length of the target}\lbl{sec:x0}
\begin{feqa}{material/HallCCryoTarget.cpp: HallCCryoTarget::targetInit}{init.f:19-44}
L_1 &= \ln 184.15 - \tfrac13\ln Z,\qquad L_2 = \ln 1194 - \tfrac23\ln Z\qquad (Z=1:\ L_1=5.31,\ L_2=6.144),\nonumber\\
f_c &= a^2\Bigl(1.202 + a^2\bigl(-1.0369 + \tfrac{1.008\,a^2}{a^2+1}\bigr)\Bigr),\quad a = Z\alpha,\qquad
X_0 = \frac{716.405\,A}{Z\,\bigl(Z(L_1-f_c)+L_2\bigr)}\ \ [\mathrm{g/cm^2}],
\lbl{eq:x0}
\end{feqa}
except $X_0 = 61.28$, 122.4 and 94.32 g/cm$^2$ for $A=1$, 2, 4 (the 1999 Particle Data Book values, as the code
comment says). The thickness $t_i$ of particle $i$ is the sum of path length over radiation length (in cm) for every
material crossed; $t_{1,2,3}$ enter the external radiation and the multiple scattering.

\subsection{Ionisation energy loss}\lbl{sec:enerloss}
For a layer of thickness $x=s\rho$ (g/cm$^2$) and a particle of energy $E$ and mass $m$ ($\gamma=E/m$,
$\beta=\sqrt{1-\gamma^{-2}}$):
\begin{feqa}{material/SimcMaterial.cpp: SimcMaterial::energyLoss}{enerloss\_new.f:19-72}
I &= 21.8\ \mathrm{eV}\ (Z=1),\quad I = 16\,Z^{0.9}\ \mathrm{eV}\ (\text{otherwise}),\nonumber\\
\hbar\omega_p &= 28.816\sqrt{\rho Z/A}\ \mathrm{eV},\qquad C_0 = \ln(\hbar\omega_p/I)+0.5,\qquad X_{10}=\log_{10}\beta\gamma,\nonumber\\
\delta_d &= \begin{cases} 0 & X_{10}<0\\
C_0 + X_{10}\ln10 + |C_0/27|\,(3-X_{10})^3 & 0\le X_{10}<3\\
C_0+X_{10}\ln10 & 3\le X_{10}<4.7\\ C_0+4.7\ln10 & X_{10}\ge4.7\end{cases}\nonumber\\
\Delta_\mathrm{mp} &= \xi\Bigl[\ln\frac{m_e}{I^2}+1.063+2\ln\beta\gamma+\ln\xi-\beta^2-\delta_d\Bigr],\qquad
\xi = 0.1536\,\frac{Z}{A}\frac{x}{\beta^2}\ [\mathrm{MeV}],\nonumber\\
\Delta E &= \Delta_\mathrm{mp} + \lambda\,\xi_L,\qquad \xi_L = \frac{0.307075}{2}\frac{Z}{A}\frac{x}{\beta^2},\qquad \lambda = -2\ln X,
\lbl{eq:enerloss}
\end{feqa}
($I$, $m_e$, $\xi$ in MeV inside the logarithms) with $X=|g|$ for a Gaussian $g$ truncated at $10\sigma$ when the loss
is sampled, $X=3$ for the minimum loss ($\lambda=-2.197$), $X=0.0067$ for the maximum ($\lambda=9.99$) and $X=1$ for
the most probable value ($\lambda=0$). A loss above the kinetic energy is limited to $E-m-10^{-7}$ MeV. The expression
has the form of the Landau most-probable loss $\xi[\ln(2m_e\beta^2\gamma^2/I)+\ln(\xi/I)+j-\beta^2-\delta]$, but with the
constant 1.063 where $\ln2+j=0.893$ for the usual $j=0.200$; the density term $\delta_d$ has no cited source in the code.
The sampled distribution of $\lambda=-2\ln|g|$ approximates the Landau tail; it is not the Landau function.

The losses enter at three places: the beam loss up to the vertex (\eqr{eq:ein}), the sampled losses $\Delta E_{2,3}$ of
the scattered particles (\code{completeEv}), and the most-probable corrections of the reconstruction (\eqr{eq:recon}).
Without \code{using\_Eloss} all are 0 and \code{correct\_Eloss} is switched off.

\subsection{Multiple scattering}
\begin{feq}{material/HallCCryoTarget.cpp: HallCCryoTarget::targetMusc}{target.f:569-575}
\theta_0^\mathrm{ms} = \frac{13.6\ \mathrm{MeV}}{p\,\beta}\sqrt{t}\,\Bigl(1+0.088\log_{10}\frac{t}{\beta^2}\Bigr),
\qquad \theta^\mathrm{ms}_{1,2} = \theta_0^\mathrm{ms}\,g_{1,2},
\lbl{eq:musc}
\end{feq}
two independent Gaussians truncated at $3.5\sigma$, added to $y'$ and $x'$; the form is that of Lynch and Dahl
\citer{lynch}, as the code comment says. The electron uses $\beta=1$.

% =====================================================================================================================
\section{Coulomb correction}\lbl{sec:coulomb}
With \code{using\_Coulomb} (switched off for $Z=1$) the electron gains, on its way in, the potential energy at the centre
of a uniformly charged sphere of charge $(Z-1)$, scaled by 0.75 following Aste et al.\ \citer{aste}:
\begin{feq}{material/HallCCryoTarget.cpp: HallCCryoTarget::targetInit}{init.f:75-78}
V_C = 0.75\cdot\frac{3}{2}\,\frac{(Z-1)\,\alpha\hbar c}{R},\qquad R = 1.1\,A^{1/3}+0.86\,A^{-1/3}\ \mathrm{fm},
\lbl{eq:coulomb}
\end{feq}
the same value for every event. The beam energy at the vertex is raised by $V_C$ (\eqr{eq:ein}), the scattered
electron loses $V_C$ on its way out (\eqr{eq:sp}), the electron-energy edges are shifted (\eqr{eq:edges}), the
reconstruction does not correct for it (\eqr{eq:einrec}), and the cross section is multiplied by the focusing factor
\begin{feq}{reactions/EventChain.cpp: assembleWeight}{event.f:1536-1538}
\sigma \leftarrow \sigma\,\bigl(1+V_C/E_b\bigr)^2.
\lbl{eq:coulombfocus}
\end{feq}
The hadron is neither decelerated nor deflected.

% =====================================================================================================================
\section{Radiative corrections}\lbl{sec:radiation}

\subsection{Model and inputs}
The radiator \code{simc\_rad} (\code{SimcRadiator}, radiation/) radiates in the peaking approximation: each charged leg
(tail 1 the incoming electron, tail 2 the scattered electron, tail 3 the outgoing hadron) emits photons collinear with
it, internal and external bremsstrahlung being combined into one power-law spectrum per tail. The formalism is that of
N.\ Makins' thesis \citer{makins}; its published description is Ent et al.\ \citer{ent}. The vertex kinematics are
those after the photon of tail 1, so the cross section is evaluated at the lowered beam energy.

\begin{simctable}{footnotesize}{L{0.2\linewidth}L{0.14\linewidth}L{0.6\linewidth}}
\caption{Inputs of the radiative corrections (one input per row).}\lbl{tab:radinputs}\\
\toprule\rowcolor{simchead}\simcth{input} & \simcth{default} & \simcth{meaning}\\\midrule\endfirsthead
\toprule\rowcolor{simchead}\simcth{input} & \simcth{default} & \simcth{meaning}\\\midrule\endhead
\bottomrule\endlastfoot
\code{using\_rad} & 0 & 1: radiate; 0: no photon and no correction factor\\
\code{rad\_flag} & 0 & 0: one tail per event, weight with the exact internal soft factor; 1: one tail per event, BASICRAD weight; 2: one tail chosen with equal probabilities; 3: all tails radiate\\
\code{extrad\_flag} & 0 & external correction $\phi$: 1 none, 2 linear, 3 Friedrich; 0 means 3 for \code{rad\_flag} 0 and 1 otherwise\\
\code{one\_tail} & 0 & restricts radiation to some tails (0 all; 1, 2, 3 one tail; $-1$, $-2$, $-3$ all but one); switching off tail 1 or 2 is not well defined (warning)\\
\code{intcor\_mode} & 0 & 0: hard correction from the Schwinger formula; otherwise from the exact soft-photon calculation\\
\code{use\_offshell\_rad} & 0 & 1: the off-shell variant of the soft-photon calculation\\
\code{use\_expon} & 0 & 1: exponentiate the soft factor\\
\code{Egamma\_gen\_max} & 0 & if $>0.01$ MeV: a fixed upper limit for every photon energy\\
\code{dE\_edge\_test} & 0 & the energy margin $d_E$ added to the photon limits\\
\end{simctable}

\subsection{Per-event constants}\lbl{sec:radconst}
At the end of \code{completeEv} \code{IRadiator::initEvent} computes the constants from the vertex kinematics. The
external radiator thicknesses use the factor $b$ (named \code{etatzai} in the code):
\begin{feqa}{radiation/RadcInit.cpp: radcInit, radcInitEv}{init.f:634, 682-688}
b &= \frac{1}{9}\Bigl(12 + \frac{Z+1}{Z L_1+L_2}\Bigr),\qquad bt_i = b\,t_i\quad(i=1,2),
\lbl{eq:bt}
\end{feqa}
($t_1$ the material before the vertex, $t_2$ after it on the electron side; the hadron has no external radiation). The
internal equivalent radiators (extended peaking approximation) are
\begin{feqa}{radiation/Radc.cpp: lambdaDave}{radc.f:792-827}
\lambda_1 &= \frac{\alpha}{\pi}\Bigl[2\ln\frac{2E_\mathrm{in}}{m_e}-1+\ln\frac{1-\cos\theta_e}{2}+2\ln\frac{E_\mathrm{in}}{E'}\cdot\Theta_3\Bigr],\nonumber\\
\lambda_2 &= \frac{\alpha}{\pi}\Bigl[2\ln\frac{2E'}{m_e}-1+\ln\frac{1-\cos\theta_e}{2}+2\ln\frac{E_\mathrm{in}}{E'}\cdot\Theta_3\Bigr],\nonumber\\
\lambda_3 &= \frac{\alpha}{\pi}\Bigl[\frac{E_h}{P_h}\ln\frac{E_h+P_h}{E_h-P_h}-2\Bigr]\cdot\Theta_3\ \ (\text{set to 0 if negative}),
\lbl{eq:lambda}
\end{feqa}
with $\Theta_3=1$ if tail 3 is enabled, 0 otherwise (the input values \code{lambda} are not used). The hadron radiates in
an event when $\lambda_3>0$. The exponents and normalisations of the power-law spectra $c\,E_\gamma^{g-1}$ are
\begin{feqa}{radiation/RadcInit.cpp: basicradInitEv}{init.f:758-810}
g_1 &= \lambda_1+bt_1,\quad g_2=\lambda_2+bt_2,\quad g_3=\lambda_3,\quad g_0=g_1+g_2+g_3,\nonumber\\
c^\mathrm{int}_{1,2} &= \frac{\lambda_{1,2}}{(E_\mathrm{in}E')^{\lambda_{1,2}/2}}\frac{e^{-\gamma_E\lambda_{1,2}}}{\Gamma(1+\lambda_{1,2})},\qquad
c^\mathrm{int}_3 = \frac{\lambda_3}{(M_pE_h)^{\lambda_3/2}}\frac{e^{-\gamma_E\lambda_3}}{\Gamma(1+\lambda_3)},\nonumber\\
c^\mathrm{ext}_{i} &= \frac{bt_{i}}{E_{i}^{bt_{i}}\,\Gamma(1+bt_{i})}\ \ (E_1=E_\mathrm{in},\ E_2=E'),\nonumber\\
c_{i} &= c^\mathrm{int}_ic^\mathrm{ext}_i\,\frac{g_i}{\lambda_i\,bt_i}\,\frac{\Gamma(1+\lambda_i)\Gamma(1+bt_i)}{\Gamma(1+g_i)}\ \ (i=1,2),\qquad c_3 = c^\mathrm{int}_3,
\lbl{eq:basicconst}
\end{feqa}
and the combined constants $c^\mathrm{int}_0$ (exponent $g_\mathrm{int}=\lambda_1+\lambda_2+\lambda_3$),
$c^\mathrm{ext}_0$ ($g_\mathrm{ext}=bt_1+bt_2$) and $c_0$ follow the convolution rule
$c_ac_b\,\frac{g_a+g_b}{g_ag_b}\frac{\Gamma(1+g_a)\Gamma(1+g_b)}{\Gamma(1+g_a+g_b)}$. The mass in $c^\mathrm{int}_3$ is $M_p$
for every hadron (a code comment notes that it should be $M_p$). The hard correction and the exponent of the combined
tail are
\begin{feqa}{radiation/RadcInit.cpp: radcInitEv, basicradInitEv}{init.f:700-715, 723-725, 809-810}
H &= \frac{1}{1-\delta_\mathrm{hard}},\qquad
g_4 = -E_c\,\frac{\mathrm{d}\delta_\mathrm{soft}}{\mathrm{d}E}\Big|_{E_c} + bt_1+bt_2,\quad E_c = 450\ \mathrm{MeV},\nonumber\\
c_4 &= \frac{g_4}{(E_\mathrm{in}E')^{g_4}\,\Gamma(1+g_4)}\ \bigl[/E_h^{g_4}\ \text{if } g_3>0\bigr],\qquad
f_i = g_i/g_0\ (i=1,2,3),
\lbl{eq:hardcor}
\end{feqa}
with $\delta_\mathrm{hard}$ and $\mathrm{d}\delta_\mathrm{soft}/\mathrm{d}E$ from the soft-photon calculation
(\secr{sec:brem}); $H$ is \code{hardcorfac} and $f_i$ are the probabilities of the tails when one tail is chosen.

\subsection{Generating the photons}\lbl{sec:basicrad}
For each radiating tail a photon energy is drawn between limits $E_\gamma^\mathrm{min}$ and $E_\gamma^\mathrm{max}$ (given
in the reaction notes; they keep the radiated event inside the edges and the energy budget) from the normalised power
law of exponent $g$:
\begin{feqa}{radiation/Radc.cpp: basicrad}{radc.f:61-84}
y &= y_\mathrm{min}+\xi(1-y_\mathrm{min}),\quad y_\mathrm{min} = (E_\gamma^\mathrm{min}/E_\gamma^\mathrm{max})^g,\quad
E_\gamma = y^{1/g}E_\gamma^\mathrm{max},\nonumber\\
w_\mathrm{basic} &= \frac{c}{g}\Bigl[(E^\mathrm{max}_\gamma)^g-(E^\mathrm{min}_\gamma)^g\Bigr],
\lbl{eq:basicrad}
\end{feqa}
where $E_\gamma^\mathrm{min}\le0$ counts as 0, the try fails if $E_\gamma^\mathrm{max}\le\max(0,E_\gamma^\mathrm{min})$, and
$g\le0$ gives no photon and $w=1$. When one tail is chosen per event (\code{rad\_flag} $\le1$) the choice uses the
probabilities $f_i$ and the photon is drawn with $(g_4,c_4)$; with \code{rad\_flag} 2 or 3 each tail uses its own
$(g_i,c_i)$. The photon of tail 1 lowers the vertex beam energy and \code{completeEv} is called again, so the vertex
kinematics, the jacobian and the radiative constants are recomputed. Tails 2 and 3 change only the \code{orig} event:
$E'^\mathrm{orig} = E' - E_{\gamma2}$ (the try fails if $\le0$) and $E_h^\mathrm{orig}=E_h-E_{\gamma3}$ (fails if
$\le M_h$).

\subsection{The radiative weight}\lbl{sec:radweight}
For \code{rad\_flag} 2 and 3 the weight is the product of $w_\mathrm{basic}\,\phi_i$ over the radiated tails, with
\begin{feqa}{radiation/Radc.cpp: extradPhi}{radc.f:692-713}
\phi^{(2)}_i &= 1-\frac{bt_i}{E_i\,g_i}E_\gamma,\qquad
\phi^{(3)}_i = \Bigl(1-x+\frac{x^2}{b}\Bigr)\exp\Bigl[\frac{bt_i}{b}\Bigl(b-\tfrac12-bx+\tfrac{x^2}{2}\Bigr)\Bigr]\Gamma(1+bt_i),
\ \ x = E_\gamma/E_i,
\lbl{eq:phi}
\end{feqa}
for \code{extrad\_flag} 2 and 3 ($\phi=1$ for \code{extrad\_flag} 1 and for tail 3). For \code{rad\_flag} 1 the weight
is $w_\mathrm{basic}$ times $1-(bt_1/E_1+bt_2/E_2)E_\gamma/(g_1+g_2)$ when \code{extrad\_flag} $\le2$. For
\code{rad\_flag} 0 (the default) the weight of the chosen tail does not depend on the drawn $E_\gamma$:
\begin{feqa}{radiation/Radc.cpp: peakedRadWeight}{radc.f:599-652}
w_\mathrm{peak} &= \frac{c^\mathrm{ext}_0}{g_\mathrm{ext}}\Bigl[e^{-\delta_\mathrm{soft}(E_\gamma^\mathrm{max})}(E_\gamma^\mathrm{max})^{g_\mathrm{ext}}
 - e^{-\delta_\mathrm{soft}(E_\gamma^\mathrm{min})}(E_\gamma^\mathrm{min})^{g_\mathrm{ext}}\Bigr]\nonumber\\
 &\quad\times\frac{e^{-\gamma_E g_4}}{\Gamma(1+g_4)}\,\frac{\Gamma(1+g_4-bt_1-bt_2)\,\Gamma(1+bt_1)\,\Gamma(1+bt_2)}{\Gamma(1+g_4)},
\lbl{eq:peaked}
\end{feqa}
with $\delta_\mathrm{soft}(E)$ the internal soft factor at the cut-off $E$ (the $E_\gamma^\mathrm{min}$ term is dropped when
$E_\gamma^\mathrm{min}\le0$; a negative weight is set to 0). It is the probability that the total radiated energy lies
between the limits, $e^{-\delta_\mathrm{soft}(E)}$ playing the role of the integrated internal spectrum. The Friedrich
external terms are computed but not used; $\Gamma(1+g_4)$ appears twice in the denominator, and with $g_4\approx0.05$
to $0.15$ the Gamma-function factor is $1+O(g_4^2)$. With \code{rad\_flag} 1 and \code{extrad\_flag} 3 the code reaches
\eqr{eq:peaked} with internal factors not set in that call (values of an earlier call). The generation weight finally
becomes
\begin{feq}{radiation/Radc.cpp: generateRad}{radc.f:521-522}
w_\mathrm{gen}\leftarrow w_\mathrm{gen}\cdot\frac{w_\mathrm{rad}}{H} = w_\mathrm{gen}\,w_\mathrm{rad}\,(1-\delta_\mathrm{hard}),
\lbl{eq:radweight}
\end{feq}
where $w_\mathrm{rad}$ is the last weight assigned (when one tail is chosen it is the peaked weight of that tail) and $H$
is the one of the last \code{completeEv} (after a tail-1 photon, at the lowered beam energy). Events in which a disabled
tail is chosen keep $w_\mathrm{rad}=1$. The tree column \code{radphot} is $E_{\gamma1}+E_{\gamma2}+E_{\gamma3}$.

\subsection{Soft and hard factors}\lbl{sec:brem}
The soft-photon calculation \code{brem} works in GeV with its own constants ($M=0.93827231$, $m_e=0.00051099906$,
$\alpha=1/137.0359895$) and rebuilds the electron angle and the recoil from \emph{elastic electron--proton kinematics}
with the same $E_\mathrm{in}$ and $E'$:
\begin{feqa}{radiation/Brem.cpp: brem}{brem.f:32-46}
k &= E_\mathrm{in}/1000,\ k'=E'/1000,\ \Delta=E_c/1000,\nonumber\\
\theta &= 2\arcsin\sqrt{\frac{M}{2k}\Bigl(\frac{k}{k'}-1\Bigr)},\qquad Q^2 = 4kk'\sin^2(\theta/2),\qquad
E_p = M+k-k',\qquad p=\sqrt{E_p^2-M^2}.
\lbl{eq:bremkin}
\end{feqa}
This is exact for H(e,e\pr p) only; for every other reaction the angle and $Q^2$ used for the soft and hard factors are
those of an elastic electron--proton event with the same energies. The soft factor is the sum of the direct and
interference terms of the electron currents (and of the proton currents when the hadron radiates):
\begin{feqa}{radiation/Brem.cpp: brem}{brem.f:64-83, 165-183}
b_{ei} &= -\frac{1}{2\pi}\ln\frac{k}{\Delta},\qquad b_{ef} = -\frac{1}{2\pi}\ln\frac{k'}{\Delta},\qquad
b_{ee} = -kk'(1-\cos\theta)\,I_{ee},\nonumber\\
\delta_\mathrm{soft} &= 2\alpha\bigl(b_{ei}+b_{ef}+b_{ee}\bigr)+[\text{proton terms}],\qquad
\delta_\mathrm{hard} = -\frac{\alpha}{\pi}\Bigl(-\frac{28}{9}+\frac{13}{6}\ln\frac{Q^2}{m_e^2}\Bigr),
\lbl{eq:brem}
\end{feqa}
where $I_{ee}$ is the integral of \code{inter} (logarithms plus Spence functions, the latter replaced by the
approximation $\mathrm{Sp}(x)=0$ for $|x|\le1$ and $-\frac12\ln^2|x|$ otherwise). The off-shell variant takes the
four-vectors of the event and places the initial proton at rest. The Schwinger form used for \code{intcor\_mode} 0 is
\begin{feqa}{radiation/Radc.cpp: schwinger}{radc.f:732-745}
\delta_\mathrm{soft}^\mathrm{S} &= \frac{\alpha}{\pi}\Bigl(\ln\frac{Q^2}{m_e^2}-1\Bigr)\ln\frac{E_\mathrm{in}E'\,b_\eta}{E_c^2},\quad
b_\eta = 1+\frac{2\nu\sin^2(\theta_e/2)}{A u},\nonumber\\
\delta_\mathrm{hard}^\mathrm{S} &= -\frac{\alpha}{\pi}\Bigl[2.166666\Bigl(\ln\frac{Q^2}{m_e^2}-1\Bigr)+\mathrm{Sp}\bigl(\cos^2\tfrac{\theta_e}{2}\bigr)-2.5893784-\tfrac12\ln^2\frac{E_\mathrm{in}}{E'}\Bigr],
\lbl{eq:schwinger}
\end{feqa}
with \code{spen} the series 27.7.2 of Abramowitz and Stegun \citer{as} truncated when a term falls below $10^{-4}$ of the
sum (at most 101 terms). In that mode the derivative of the soft factor is not computed and $g_4$ uses a value left
from earlier.

% =====================================================================================================================
\section{Common cross-section ingredients}\lbl{sec:xsecbase}
\subsection{Mott cross section}
\begin{feq}{physics/proton/ProtonCrossSections.cpp: sigMott}{physics\_proton.f:186-187}
\sigma_\mathrm{Mott}(E',\theta,Q^2) = \Bigl(\frac{2\alpha\hbar c\,E'\cos(\theta/2)}{Q^2}\Bigr)^2\times10^4\quad[\mu\mathrm{b/sr}],
\lbl{eq:mott}
\end{feq}
the factor $10^4$ converting fm$^2$ to $\mu$b. With $Q^2 = 4E_\mathrm{in}E'\sin^2(\theta/2)$ it becomes
$\sigma_\mathrm{Mott} = \alpha^2(\hbar c)^2\cos^2(\theta/2)\,/\,4E_\mathrm{in}^2\sin^4(\theta/2)$, the Mott cross
section of a point charge for the incident energy, without a recoil factor.

\subsection{Proton form factors}\lbl{sec:ff}
The proton form factors are the fit of Bosted \citer{bosted} (``Eqs.~4 and 5'' according to the code comment), as a
function of $Q=\sqrt{Q^2}$ in GeV:
\begin{feqa}{physics/proton/ProtonCrossSections.cpp: fofa\_best\_fit}{physics\_proton.f:150-163}
G_E &= \frac{1}{1+0.62Q+0.68Q^2+2.8Q^3+0.83Q^4},\nonumber\\
G_M &= \frac{\mu_p}{1+0.35Q+2.44Q^2+0.5Q^3+1.04Q^4+0.34Q^5},\qquad \mu_p = 2.793,
\lbl{eq:bosted}
\end{feqa}
with $Q^2$ in GeV$^2$ ($Q^2$[MeV$^2$]$\times10^{-6}$; the routine receives $-Q^2/(\hbar c)^2$ and converts it back) and
$Q=\sqrt{\max(Q^2,0)}$. The alternative ``Eqs.~6'' (one common denominator) is present as a comment and not used. The
coefficients are transcribed from the code; they have not been compared with the paper for these notes.

\begin{simctable}{footnotesize}{L{0.18\linewidth}L{0.12\linewidth}L{0.12\linewidth}L{0.12\linewidth}L{0.12\linewidth}L{0.12\linewidth}}
\caption{Coefficients of the denominators of \eqr{eq:bosted} (one form factor per row).}\lbl{tab:bosted}\\
\toprule\rowcolor{simchead}\simcth{} & \simcth{$Q$} & \simcth{$Q^2$} & \simcth{$Q^3$} & \simcth{$Q^4$} & \simcth{$Q^5$}\\\midrule\endfirsthead
\bottomrule\endlastfoot
$G_E$ & 0.62 & 0.68 & 2.8 & 0.83 & --\\
$G_M/\mu_p$ & 0.35 & 2.44 & 0.5 & 1.04 & 0.34\\
\end{simctable}

% =====================================================================================================================
\section{Validation of the common parts}\lbl{sec:validation}
The common routines are tested routine by routine against tables computed with the Fortran (CTest label \code{l1}, with
the tolerance 0 in compat mode unless the test says otherwise), and as a whole by the reaction runs of the reaction
notes. \tabr{tab:tests} lists the tests of the common parts.

\begin{simctable}{footnotesize}{L{0.36\linewidth}L{0.14\linewidth}L{0.44\linewidth}}
\caption{CTest tests of the common parts (one test per row; run with \code{ctest -R <name>}).}\lbl{tab:tests}\\
\toprule\rowcolor{simchead}\simcth{test} & \simcth{labels} & \simcth{what it checks}\\\midrule\endfirsthead
\toprule\rowcolor{simchead}\simcth{test} & \simcth{labels} & \simcth{what it checks}\\\midrule\endhead
\bottomrule\endlastfoot
\code{l1\_physics\_angles} & l1 & \eqr{eq:physang} against the Fortran table\\
\code{l1\_spectrometer\_angles} & l1 & \eqr{eq:specang}\\
\code{l1\_gauss1\_fixedrng} & l1 & truncated Gaussian numbers\\
\code{l1\_target\_init} & l1 & $X_0$, average losses, $V_C$\\
\code{l1\_extreme\_trip\_thru\_target} & l1 & extreme losses and multiple-scattering angles\\
\code{l1\_trip\_thru\_target\_fixedrng} & l1 & target and window paths and losses\\
\code{l1\_enerloss\_new} & l1 & \eqr{eq:enerloss}\\
\code{l1\_target\_musc\_fixedrng} & l1 & \eqr{eq:musc}\\
\code{l1\_lambda\_dave} & l1 & \eqr{eq:lambda}\\
\code{l1\_radc\_init\_ev} & l1 & per-event radiative constants\\
\code{l1\_basicrad\_init\_ev} & l1 & \eqr{eq:basicconst}\\
\code{l1\_basicrad\_fixedrng} & l1 & \eqr{eq:basicrad}\\
\code{l1\_peaked\_rad\_weight} & l1 & \eqr{eq:peaked}\\
\code{l1\_extrad\_phi} & l1 & \eqr{eq:phi}\\
\code{l1\_generate\_rad\_fixedrng} & l1 & photon limits and the radiative weight\\
\code{l1\_brem} & l1 & \eqr{eq:brem}\\
\code{l1\_bremos} & l1 & the off-shell variant\\
\code{l1\_schwinger} & l1 & \eqr{eq:schwinger}\\
\code{l1\_sigmott} & l1 & \eqr{eq:mott}\\
\code{l1\_fofa\_best\_fit} & l1 & \eqr{eq:bosted}\\
\code{fortran\_diff\_D006\_pass\_cuts} & fortran\_diff & \code{pass\_cuts} in both modes\\
\code{fortran\_diff\_D028\_spectrometer\_angles\_nan} & fortran\_diff & the NaN slope of \eqr{eq:specang} in both modes\\
\code{fortran\_diff\_D017\_extreme\_trip\_hadron\_angle} & fortran\_diff & the hadron extreme path in both modes\\
\end{simctable}

% =====================================================================================================================
\section{Fortran origin and Fortran findings}\lbl{sec:origin}
\begin{simctable}{footnotesize}{L{0.42\linewidth}L{0.55\linewidth}}
\caption{Fortran routine $\to$ C++ function, common parts (one routine per row).}\lbl{tab:map}\\
\toprule\rowcolor{simchead}\simcth{Fortran (upstream 59ccdf6)} & \simcth{C++ (src/...)}\\\midrule\endfirsthead
\toprule\rowcolor{simchead}\simcth{Fortran (upstream 59ccdf6)} & \simcth{C++ (src/...)}\\\midrule\endhead
\bottomrule\endlastfoot
\file{simc.f} main loop, l.176-358 & \csrc{driver/SimcRun.cpp: SimcRun::step}\\
\file{simc.f} pass\_cuts, l.236-253 & \csrc{driver/SimcRun.cpp: PassCuts}\\
\file{simc.f} targetfac and luminosity, l.101-108 & \csrc{setup/SetupBuilder.cpp: BuildRunSetup}\\
\file{simc.f} normfac and genvol, l.375-405 & \csrc{driver/RunReport.cpp: Normalise}\\
\file{simc.f} calculate\_central & \csrc{driver/Central.cpp: calculateCentral}\\
\file{simc.f} montecarlo & \csrc{driver/MonteCarlo.cpp: MonteCarlo::run}\\
\file{event.f} generate, beam part l.147-212 & \csrc{reactions/StandardBeam.cpp: StandardBeam::generate}\\
\file{event.f} generate, reaction part & \csrc{IReaction::generate} of each reaction\\
\file{event.f} complete\_ev & \csrc{reactions/EepReactionBase.cpp: EepReactionBase::completeEv}\\
\file{event.f} complete\_recon\_ev & \csrc{reactions/EepReactionBase.cpp: EepReactionBase::completeRecon}\\
\file{event.f} complete\_main & \csrc{reactions/EepReactionBase.cpp: EepReactionBase::completeMain}\\
\file{event.f} complete\_main, weight part & \csrc{reactions/EventChain.cpp: assembleWeight}\\
\file{event.f} limits\_update & \csrc{reactions/EventChain.cpp: limitsUpdate}\\
\file{event.f} physics\_angles & \csrc{setup/KinematicsUtils.cpp: PhysicsAngles}\\
\file{event.f} spectrometer\_angles & \csrc{setup/KinematicsUtils.cpp: SpectrometerAngles}\\
\file{init.f} limits\_init & \csrc{setup/SetupLimits.cpp: BuildLimits}\\
\file{init.f} target\_init & \csrc{material/HallCCryoTarget.cpp: HallCCryoTarget::targetInit}\\
\file{init.f} radc\_init & \csrc{radiation/RadcInit.cpp: radcInit}\\
\file{init.f} radc\_init\_ev & \csrc{radiation/RadcInit.cpp: radcInitEv}\\
\file{init.f} basicrad\_init\_ev & \csrc{radiation/RadcInit.cpp: basicradInitEv}\\
\file{target.f} trip\_thru\_target & \csrc{material/HallCCryoTarget.cpp: HallCCryoTarget::trip}\\
\file{target.f} extreme\_trip\_thru\_target & \csrc{material/HallCCryoTarget.cpp: extremeTripThruTarget}\\
\file{target.f} target\_musc & \csrc{material/HallCCryoTarget.cpp: targetMusc}\\
\file{enerloss\_new.f} & \csrc{material/SimcMaterial.cpp: SimcMaterial::energyLoss}\\
\file{gauss1.f} & \csrc{core/random/Gauss1.cpp: gauss1Polar}\\
call\_ranlux.f, cern/ranlux.f & \csrc{core/random/RanluxRandom.cpp: RanluxRandom}\\
\file{radc.f} basicrad & \csrc{radiation/Radc.cpp: basicrad}\\
\file{radc.f} generate\_rad & \csrc{radiation/Radc.cpp: generateRad}\\
\file{radc.f} peaked\_rad\_weight & \csrc{radiation/Radc.cpp: peakedRadWeight}\\
\file{radc.f} extrad\_friedrich & \csrc{radiation/Radc.cpp: extradFriedrich}\\
\file{radc.f} extrad\_phi & \csrc{radiation/Radc.cpp: extradPhi}\\
\file{radc.f} schwinger & \csrc{radiation/Radc.cpp: schwinger}\\
\file{radc.f} spen & \csrc{radiation/Radc.cpp: spen}\\
\file{radc.f} lambda\_dave & \csrc{radiation/Radc.cpp: lambdaDave}\\
\file{brem.f} brem & \csrc{radiation/Brem.cpp: brem}\\
\file{brem.f} bremos & \csrc{radiation/Brem.cpp: bremos}\\
\file{brem.f} inter & \csrc{radiation/Brem.cpp: inter}\\
\file{brem.f} spence & \csrc{radiation/Brem.cpp: spence}\\
\file{physics\_proton.f} sigMott & \csrc{physics/proton/ProtonCrossSections.cpp: sigMott}\\
\file{physics\_proton.f} fofa\_best\_fit & \csrc{physics/proton/ProtonCrossSections.cpp: fofa\_best\_fit}\\
\file{dbase.f} reaction flags & \csrc{core/RunSetup.cpp: DeriveReactionFlags}\\
\file{dbase.f} masses and kinematics & \csrc{setup/SetupReaction.cpp: BuildReactionAndKinematics}\\
\end{simctable}

\begin{simcfindings}{footnotesize}{L{0.19\linewidth}L{0.47\linewidth}L{0.3\linewidth}}
\caption{Fortran findings that concern every reaction (FORTRAN\_DIFFERENCES.md; one finding per row).}\lbl{tab:findings}\\
\toprule\rowcolor{simchead}\simcth{id} & \simcth{Fortran behaviour} & \simcth{default / compat}\\\midrule\endfirsthead
\toprule\rowcolor{simchead}\simcth{id} & \simcth{Fortran behaviour} & \simcth{default / compat}\\\midrule\endhead
\bottomrule\endlastfoot
\finding{D-006 pass\_cuts} & the hadron $\delta$ upper cut uses the electron limit (\file{simc.f:244}) & hadron limit / Fortran\\
\finding{D-028 F2} & NaN $y'$ from \code{spectrometer\_angles} (\file{event.f:1645}) & $y'=0$ / NaN\\
\finding{D-017} & the hadron extreme path uses the electron angle (\file{target.f:501}) & hadron angle / Fortran\\
\finding{D-024 grnd dEbeam=0} & a random number is drawn for $E_\mathrm{in}$ even when $\Delta E_b=0$ & kept in both\\
\finding{D-024 Eloss gaussians} & energy-loss Gaussians drawn without \code{using\_Eloss} & kept in both\\
\finding{D-024 stale orig} & a failed radiation leaves \code{orig} of the previous try & kept in both\\
\finding{D-019 rad\_weight} & the peaked weight overwrites the product of the radiated tails & kept in both\\
\finding{D-019 tail-1 COMMON} & $H$ recomputed by the second \code{completeEv} after a tail-1 photon & kept in both\\
\finding{D-019 intcor\_mode 0} & $g_4$ without the soft-factor derivative & kept in both\\
\finding{D-019 Q4 phi\_ext} & the external terms of \eqr{eq:peaked} unused & kept in both\\
\finding{D-019 Q4 dsoft\_int} & \code{rad\_flag} 1 with \code{extrad\_flag} 3 uses stale internal factors & kept in both\\
\finding{D-019 one\_tail} & a disabled tail can be chosen (no photon) & kept in both\\
\finding{D-019 doing\_proton} & the tail-3 switch also removes the interference term of $\lambda_{1,2}$ & kept in both\\
\finding{D-005 lambda input} & the input $\lambda$ values are overwritten & kept in both\\
\finding{P5.F1 Q8} & the soft and hard factors use elastic e--p kinematics & kept in both\\
\finding{D-012 F2} & guards of the Gaussian and of the logarithm in the energy loss & kept in both\\
\finding{D-012 F3} & energy loss limited to the kinetic energy & kept in both\\
\finding{D-012 F4} & no windows for the calorimeter arms & kept in both\\
\finding{D-012 F5} & Gaussian truncations $3.5\sigma$, $10\sigma$ & kept in both\\
\finding{D-019 gamma} & the polynomial \code{gamma} of \file{radc.f} is not used (library $\Gamma$) & as the Fortran binary\\
\finding{P8.2 targetfac} & quadruple-precision \code{targetfac} & reproduced\\
\finding{D-002 RNG} & RANLUX in the Fortran & \code{mt64} by default, \code{ranlux} on request\\
\end{simcfindings}

\begin{simcrefs}
\simcref{ent} R.~Ent et al., ``Radiative corrections for (e,e\pr p) reactions at GeV energies'', Phys.\ Rev.\ C
  \textbf{64}, 054610 (2001). Not cited in the code; the published account of the formalism.
\simcref{makins} N.~C.~R.~Makins, Ph.D.\ thesis, MIT (1994), Sect.~4.5.6 (cited at \fsrc{radc.f:161}).
\simcref{lynch} G.~R.~Lynch and O.~I.~Dahl, Nucl.\ Instrum.\ Meth.\ B \textbf{58}, 6 (1991), Eq.~6 (cited at
  \fsrc{target.f:568}).
\simcref{aste} A.~Aste, K.~Hencken, J.~Jourdan, I.~Sick, D.~Trautmann, Eur.\ Phys.\ J.\ A \textbf{26}, 167 (2005)
  (cited at \fsrc{init.f:71}).
\simcref{mo} L.~W.~Mo and Y.~S.~Tsai, Rev.\ Mod.\ Phys.\ \textbf{41}, 205 (1969).
\simcref{as} M.~Abramowitz and I.~A.~Stegun, \emph{Handbook of Mathematical Functions}, formula 27.7.2 (cited at
  \fsrc{radc.f:759-760}).
\simcref{bosted} P.~E.~Bosted, ``Empirical fit to the nucleon electromagnetic form factors'', Phys.\ Rev.\ C
  \textbf{51}, 409 (1995), Eqs.~4 and 5 (cited at \fsrc{physics\_proton.f:139-141}).
\end{simcrefs}
\end{document}
